\documentclass[12pt,a4paper]{article}
\usepackage[utf8]{inputenc}
\usepackage[margin=1in]{geometry}
\usepackage{setspace}
\usepackage{cite}
\usepackage{hyperref}
\usepackage{booktabs}
\usepackage{longtable}

\title{Indian Supreme Court Defines Hierarchical Classification Framework for Food Products, Overruling Common Parlance Precedents}

\author{Lalitha A R\\
\texttt{lalithaar.research@gmail.com}\\
Lead Researcher\\
Interdisciplinary Systems Research Lab (ISRL)}

\date{\today}

\begin{document}

\maketitle

\begin{abstract}
This report synthesizes landmark judicial decisions in Indian food classification law, documenting the transition from common parlance-based interpretation to a hierarchical technical classification framework. The January 6, 2026 Supreme Court judgment in \textit{Commissioner of Customs Import v. Ms Welkin Foods} formally established a precedential hierarchy that prioritizes statutory interpretation over lay understanding, marking a departure from pre-HSN era jurisprudence. This analysis examines four domains (GST/Tax, Food Safety, Customs, and Dietary Labels) where technical definitions now supersede common parlance, with significant implications for food informatics, regulatory compliance, and ingredient classification systems.
\end{abstract}

\section{Introduction}

The classification of food products under Indian law has undergone a fundamental transformation. Historically, courts relied heavily on common parlance---the everyday understanding of terms as used in ordinary commerce---to interpret ambiguous statutory language. This approach, exemplified by pre-1986 cases such as \textit{Nix v. Hedden} (1893) in the United States and \textit{Ramavatar Budhaiprasad v. Assistant Sales Tax Officer} (1961) in India, prioritized accessibility and commercial understanding over technical precision.

The adoption of the Harmonised System of Nomenclature (HSN) by India in 1986, following its international introduction in 1988 by the World Customs Organisation, initiated a gradual shift toward technical classification. However, the tension between common understanding and scientific taxonomy persisted until the Supreme Court's 2026 ruling explicitly established a hierarchy for resolving classification disputes.

This report documents this watershed moment and its implications across multiple regulatory domains.

\section{Historical Context: Pre-HSN Era Common Parlance Precedents}

\subsection{United States: \textit{Nix v. Hedden} (1893)}

The seminal case \textit{Nix v. Hedden}, 149 U.S. 304 (1893), established that tomatoes should be classified as vegetables rather than fruits for tariff purposes, despite their botanical classification. The Supreme Court of the United States held that classification should follow common parlance---how ordinary people in commerce understand terms---rather than technical botanical definitions.

\subsection{India: Early Common Parlance Cases}

\subsubsection{\textit{Ramavatar Budhaiprasad v. Assistant Sales Tax Officer} (1961)}

In this landmark case, AIR 1961 SC 1325; (1962) 1 SCR 279; (1961) 12 STC 286, decided on March 14, 1961, the Supreme Court of India addressed whether betel leaves should be classified as vegetables for sales tax exemption purposes under the Central Provinces and Berar Sales Tax Act, 1947.

The Court held that the word ``vegetables'' must be interpreted not in a technical or botanical sense, but in its popular sense as understood in common language---denoting classes of vegetable matter grown in kitchen gardens or farms and used for the table. The Court stated: ``It has not been defined in the Act and being a word of everyday use it must be construed in its popular sense, meaning that sense which people conversant with the subject matter with which the statute is dealing would attribute to it.''

The Court ruled that betel leaves, while botanically plant matter, were not vegetables in common parlance and were therefore taxable.

\subsubsection{\textit{Krishna Iyer v. State of Kerala} (1962)}

Decided on March 6, 1962, this case from the Kerala High Court similarly applied the common parlance test to determine whether green ginger qualified as a vegetable for tax exemption purposes. The Court held that vegetables should be understood ``as commonly understood denoting those classes of vegetable matter which are grown in kitchen gardens and are used for the table,'' and concluded that green ginger, despite being plant matter, was included in the specific term ``ginger'' in the tax schedule and was therefore taxable.

\subsection{The Pre-HSN Framework}

Prior to India's adoption of the HSN system in 1986, courts consistently applied common parlance as the primary interpretive tool for ambiguous statutory terms. This approach served several purposes:

\begin{itemize}
\item It made tax classifications accessible to ordinary merchants without specialized knowledge
\item It aligned legal interpretations with commercial practice
\item It avoided the complexity of technical botanical or chemical classifications
\item It provided predictability based on everyday understanding
\end{itemize}

\section{The Watershed Moment: \textit{Commissioner of Customs Import v. Ms Welkin Foods} (2026)}

\subsection{Case Details}

\textbf{Citation:} 2026 SCC OnLine SC 27; 2026 INSC 19\\
\textbf{Date:} January 6, 2026 (reported January 6-7, 2026)\\
\textbf{Court:} Supreme Court of India\\
\textbf{Bench:} Justice J.B. Pardiwala and Justice R. Mahadevan\\
\textbf{Parties:} Commissioner of Customs (Import) v. M/s Welkin Foods

\subsection{Facts and Issue}

The case concerned the proper classification of imported aluminium shelving used for mushroom cultivation. The respondent, Welkin Foods, argued the goods should be classified under Customs Tariff Item (CTI) 84369900 as ``parts'' of agricultural machinery. The Revenue contended the shelving should be classified under CTI 76109010 as ``Aluminium Structures.''

The core legal question was whether the intended use of the product (mushroom cultivation) or its objective technical characteristics should govern classification.

\subsection{The Court's Reasoning}

The Supreme Court held that classification must be based on objective characteristics of the product, not solely on intended end-use. The Court established several critical principles:

\begin{enumerate}
\item \textbf{Structure vs. Machine:} The shelving was held to be a ``structure'' (fixed in place) rather than a ``part'' of a machine. It did not qualify as a component essential for the mechanical function of agricultural machinery.

\item \textbf{Material Identity Primacy:} While exclusive use can sometimes influence classification, it does not override the fundamental material identity of the product when it is specifically described elsewhere in the tariff.

\item \textbf{Hierarchical Framework:} Most significantly, the Court articulated a hierarchy for classification disputes, stating: ``It is only in a state of statutory silence, where the legislative intent remains unexpressed, that the tribunals or courts may resort to the common or trade parlance test.''
\end{enumerate}

\subsection{The Established Hierarchy}

The \textit{Ms Welkin Foods} judgment established the following precedential hierarchy for food and product classification:

\begin{enumerate}
\item \textbf{Judicial Interpretation of Statute} (highest priority)---How the court reads the HSN and statutory provisions
\item \textbf{Technical/Scientific Definition}---When statute provides technical guidance through HSN codes
\item \textbf{Expert Opinion}---Testimony from qualified experts in relevant fields
\item \textbf{Common Parlance}---Trade usage and ordinary understanding (only in statutory silence)
\end{enumerate}

This hierarchy fundamentally reorients Indian classification jurisprudence, relegating common parlance from its historical primacy to a fallback position.

\section{Domain-Specific Applications}

The hierarchical framework established in \textit{Ms Welkin Foods} has been applied consistently across multiple regulatory domains, demonstrating the pervasiveness of technical classification over common understanding.

\subsection{GST and Tax Domain: Scientific Composition Prevails}

\subsubsection{\textit{In re Gajanand Foods Private Limited}}

The Gujarat Authority for Advance Ruling (GAAR) and subsequently the Appellate Authority (AAAR) addressed whether instant mix flours containing spices, leavening agents, and other additives should be classified under Chapter Headings 1102 or 1106 (attracting 5\% GST) or under Heading 2106 90 (attracting 18\% GST).

\textbf{Ruling:} The AAAR held that instant mix flours for products like Gota, Khaman, Dhokla, Idli, Dosa, Handvo, and others, containing 5-37\% additional ingredients (spices, salt, sodium bicarbonate, chili powder), are classifiable under HSN 2106 90 as ``Food Preparations not elsewhere specified or included,'' attracting 18\% GST.

\textbf{Rationale:} The technical composition, including functional additives that transformed the product from mere flour into a meal preparation kit, removed it from the common category of ``flour.'' The presence of leavening agents, spices, and other ingredients meant for creating a specific dish demonstrated that these were food preparations, not basic flours.

\subsubsection{\textit{In re Ramdev Food Products Private Limited}}

In a parallel case, the Gujarat AAAR addressed similar instant mix flours produced by Ramdev Food Products, including instant mixes for Gota, Khaman, Dalwada, Dahiwada, Idli, Dhokla, Dosa, Pizza, Methi Gota, and Handvo.

\textbf{Ruling:} The AAAR upheld the AAR's classification of these products under HSN 2106 90, attracting 18\% GST. The Court rejected arguments based on VAT-era precedents, holding that ``merely because the end consumer of the Instant Mix Flour is required to follow certain food preparation processes before such product(s) can be consumed, is no ground to take these products out of Chapter Heading 2106.''

\textbf{Key Principle:} The technical composition and processing state---not the trade name or common understanding---governs classification under the HSN system.

\subsection{Food Safety Domain: Nutrient Thresholds Over Marketing Terms}

\subsubsection{\textit{3S and Our Health Society v. Union of India}}

\textbf{Case Details:} Writ Petition (Civil) No. 437/2024\\
\textbf{Court:} Supreme Court of India\\
\textbf{Bench:} Justice J.B. Pardiwala and Justice R. Mahadevan (initial disposal: April 9, 2025)\\
\textbf{Subsequent hearings:} February 2026

This ongoing public interest litigation seeks mandatory Front-of-Package Warning Labels (FoPWL) on packaged food products containing high levels of sugar, salt, and saturated fats.

\textbf{Court's Direction:} The Supreme Court directed the Food Safety and Standards Authority of India (FSSAI) to prioritize scientific thresholds of salt, sugar, and saturated fats over the food industry's preferred marketing terminology. The Court emphasized that consumer health protection requires objective, scientifically determined nutrient levels rather than subjective or trade-based descriptions.

\textbf{Implication:} The Court's insistence on scientific measurement over industry terminology reflects the same hierarchical principle established in \textit{Ms Welkin Foods}---technical, objective criteria supersede commercial nomenclature.

\subsection{Customs Domain: Engineering Function Over Trade Nomenclature}

The \textit{Ms Welkin Foods} case itself exemplifies this domain. The Supreme Court held that aluminium racks for mushroom cultivation are technically structures (Chapter 76) and not machinery (Chapter 84) because they lack mechanical function, regardless of their trade name or intended agricultural use.

This establishes that in customs classification, the technical characteristics---material composition and functional properties---override the commercial designation or end-use of a product.

\subsection{Dietary Labels Domain: Biological Origin Disclosure Mandatory}

\subsubsection{\textit{Ram Gaua Raksha Dal v. Union of India and Others}}

\textbf{Case Details:} W.P.(C) 12055/2021\\
\textbf{Court:} Delhi High Court\\
\textbf{Bench:} Justice Vipin Sanghi and Justice Jasmeet Singh (December 2021) / Justice Vipin Sanghi and Justice Dinesh Kumar Sharma (March 2022)\\
\textbf{Date:} December 9, 2021; subsequent order March 2, 2022

This case challenged the inadequate disclosure of animal-sourced ingredients in packaged food products, particularly where International Numbering System (INS) codes obscure the biological origin of food additives.

\textbf{Court's Ruling:} The Delhi High Court held that the biological origin of ingredients must be disclosed, stating that ``every person has a right to know as to what he/she is consuming, and nothing can be offered to the person on a platter by resort to deception, or camouflage.''

The Court directed that:

\begin{itemize}
\item Food Business Operators must make full and complete disclosure of all ingredients, not only by their code names (INS numbers) but also by disclosing whether they originate from plant, animal source, or are manufactured in a laboratory.
\item The disclosure must specify the actual plant or animal source, regardless of the percentage used in the food article.
\item Even minuscule amounts of animal-sourced ingredients (other than milk, milk products, honey, beeswax, carnauba wax, or shellac) render the product non-vegetarian and must be disclosed accordingly.
\item Chemical code alone ``camouflages the truth from the consumer.''
\end{itemize}

\textbf{Constitutional Basis:} The Court grounded this requirement in Articles 19(1)(a) (freedom of speech and information), 21 (right to life and health), and 25 (freedom of religion) of the Indian Constitution, recognizing that dietary choices based on religious, ethical, or health considerations require transparent ingredient disclosure.

\textbf{Implication:} This case demonstrates that in labeling disputes, the biological or chemical origin---a technical, scientific classification---supersedes simplified marketing designations or chemical code nomenclature.

\section{Synthesis: The Four-Domain Framework}

Table \ref{tab:domains} synthesizes the current state of classification law across the four primary domains:

\begin{table}[h]
\centering
\caption{Classification Framework Across Regulatory Domains}
\label{tab:domains}
\begin{tabular}{@{}lll@{}}
\toprule
\textbf{Domain} & \textbf{Classification Winner} & \textbf{Landmark Case} \\ \midrule
GST/Tax & Technical (HSN) & \textit{Gajanand Foods}; \textit{Ramdev Food Products} \\
Food Safety & Scientific (Nutrient Levels) & \textit{3S and Our Health Society} \\
Customs & Technical (Engineering) & \textit{Ms Welkin Foods} \\
Dietary Labels & Biological Origin & \textit{Ram Gaua Raksha Dal} \\ \bottomrule
\end{tabular}
\end{table}

\section{Mandatory Disclosure Requirements: Implications for Food Informatics}

The cases analyzed in this report establish several mandatory disclosure requirements under Indian law. These requirements have direct implications for food informatics systems, ingredient databases, and regulatory compliance frameworks.

\subsection{Source Disclosure}

\textbf{Requirement:} Even if an ingredient is heavily processed or becomes a derivative compound, the source must be declared.

\textbf{Example:} Lecithin must be labeled as ``Lecithin (Soy)'' or ``Lecithin (Egg),'' not merely as ``Lecithin'' or by INS code alone.

\textbf{Legal Basis:} \textit{Ram Gaua Raksha Dal v. Union of India}

\subsection{Dietary Status}

\textbf{Requirement:} Products must be classified as vegetarian (with egg/dairy), non-vegetarian, or pure vegetarian (no animal source or dairy).

\textbf{Principle:} Even if an ingredient is a chemical derived from an animal source, it must be declared with respect to its source.

\textbf{Legal Basis:} \textit{Ram Gaua Raksha Dal v. Union of India}; Food Safety and Standards (Labelling and Display) Regulations, 2020

\subsection{Allergen Declaration}

\textbf{Requirement:} Even if processing is extreme and the allergen potency is reduced from the source level, allergen presence must still be declared for consumer safety.

\textbf{Legal Basis:} Right to health (Article 21 of the Constitution); FSSAI allergen declaration requirements

\subsection{Functional Class for Chemicals}

\textbf{Requirement:} For chemical additives, the functional class (purpose of inclusion) must be declared followed by the INS Number.

\textbf{Example:} ``Preservative (INS 202)'' rather than merely ``INS 202'' or ``Potassium Sorbate.''

\textbf{Legal Basis:} FSSAI Labelling Regulations 2020; functional usage declaration requirements

\section{Analytical Implications for Food Informatics}

The hierarchical framework and mandatory disclosure requirements have profound implications for food informatics systems:

\subsection{Attribute-Based Ingredient Classification}

The cases reveal that the determination of whether something constitutes a separate ingredient entity versus a variant of an existing entity depends on functional attributes rather than source alone.

\textbf{Key Principle:} Functionality takes precedence over source material when determining ingredient separateness.

For example, in the instant mix flour cases, the presence of functional additives (leavening agents, spices used for specific culinary purposes) transformed what might be considered a ``variant of flour'' into a distinct food preparation. The additives were not merely processing aids but functional components that changed the nature of the product.

\subsection{Hierarchical Data Modeling}

Food informatics databases must now implement hierarchical classification systems that mirror the judicial hierarchy:

\begin{enumerate}
\item \textbf{Statutory Classification Layer:} HSN codes and tariff classifications
\item \textbf{Technical/Scientific Layer:} Chemical composition, functional properties, biological origin
\item \textbf{Commercial Layer:} Trade names, common parlance terms (lowest priority)
\end{enumerate}

\subsection{Mandatory Metadata Requirements}

Any comprehensive food ingredient database must now capture:

\begin{itemize}
\item Biological/chemical source (plant, animal, synthetic)
\item Specific source species/material (even if heavily processed)
\item Functional class (for additives)
\item Allergen status (regardless of processing degree)
\item Dietary classification (vegetarian, vegan, non-vegetarian)
\item HSN code classification
\end{itemize}

\section{Conclusion}

The January 6, 2026 Supreme Court judgment in \textit{Commissioner of Customs Import v. Ms Welkin Foods} represents a watershed moment in Indian food classification jurisprudence. By formally establishing a hierarchical framework that prioritizes statutory interpretation and technical definition over common parlance, the Court has fundamentally reoriented how food products are classified across multiple regulatory domains.

This shift reflects the increasing complexity of food systems, the globalization of trade through standardized systems like HSN, and the constitutional imperative for transparent disclosure that enables informed consumer choice. The historical reliance on common parlance, while accessible and commercially grounded, proved insufficient in an era of complex food processing, international trade nomenclature, and diverse dietary requirements based on health, religion, and ethics.

The four-domain analysis presented in this report---spanning GST/Tax, Food Safety, Customs, and Dietary Labels---demonstrates the consistency with which Indian courts are now applying technical classification principles. In each domain, technical or scientific attributes supersede lay understanding or commercial nomenclature.

For food informatics systems, regulatory compliance frameworks, and ingredient databases, these decisions mandate a fundamental restructuring. Classification systems must be hierarchical, metadata must capture technical attributes (source, function, composition), and disclosure must prioritize scientific accuracy over commercial simplicity.

This report serves as a synthesis of century-spanning jurisprudential evolution, documenting the transition from common parlance dominance to technical hierarchy supremacy. It provides legal researchers, food \& beverage lawyers, compliance professionals, and informatics specialists with a comprehensive framework for understanding and applying current Indian food classification law.

\section*{Acknowledgments}

My deepest gratitude to Mr. Krishna, whose constancy forms the foundation upon which all my work, including this, quietly rests.

Salutations to the Goddess who dwells in all beings in the form of intelligence. I bow to her again and again.

This report was prepared as part of the Indian Food Informatics Data (IFID) project at the Interdisciplinary Systems Research Lab (ISRL). The synthesis draws upon extensive legal research and domain analysis conducted for food informatics applications.

\begin{thebibliography}{99}

\bibitem{nix1893}
\textit{Nix v. Hedden}, 149 U.S. 304 (1893), Supreme Court of the United States.

\bibitem{ramavatar1961}
\textit{Ramavatar Budhaiprasad v. Assistant Sales Tax Officer, Akola}, AIR 1961 SC 1325; (1962) 1 SCR 279; (1961) 12 STC 286, Supreme Court of India (March 14, 1961).

\bibitem{krishnaiyer1962}
\textit{K. Krishna Iyer v. State of Kerala}, Kerala High Court (March 6, 1962).

\bibitem{welkinfoods2026}
\textit{Commissioner of Customs (Import) v. M/s Welkin Foods}, 2026 SCC OnLine SC 27; 2026 INSC 19, Supreme Court of India (January 6, 2026).

\bibitem{gajanandfoods}
\textit{In re Gajanand Foods Private Limited}, Advance Ruling No. GUJ/GAAR/R/28/2021 (July 19, 2021); Gujarat Appellate Authority for Advance Ruling (2024).

\bibitem{ramdevfoods}
\textit{In re Ramdev Food Products Private Limited}, Gujarat Authority for Advance Ruling (July 19, 2021); Gujarat Appellate Authority for Advance Ruling (2025).

\bibitem{3shealth}
\textit{3S and Our Health Society v. Union of India and Others}, Writ Petition (Civil) No. 437/2024, Supreme Court of India (April 9, 2025; subsequent hearings February 2026).

\bibitem{ramgau}
\textit{Ram Gaua Raksha Dal v. Union of India and Others}, W.P.(C) 12055/2021, Delhi High Court (December 9, 2021; March 2, 2022).

\bibitem{fssai2020}
Food Safety and Standards Authority of India, Food Safety and Standards (Labelling and Display) Regulations, 2020.

\bibitem{customstariff}
Customs Tariff Act, 1975 (51 of 1975), Government of India.

\bibitem{hsn}
World Customs Organisation, Harmonised System of Nomenclature (1988); adopted by India (1986).

\bibitem{constitution}
Constitution of India, 1950, Articles 19(1)(a), 21, 25.

\end{thebibliography}

\end{document}